\begin{thebibliography}{23}

\bibitem[1]{ref1}
AmirMahdi Ahmadinejad, Jonathan Kelner, Jack Murtagh, John Peebles, Aaron Sidford, and Salil Vadhan. High-precision estimation of random walks in small space. In Proceedings of the 61st Annual IEEE Symposium on Foundations of Computer Science (FOCS 2020), pages 1295–1306. IEEE, 2020.

\bibitem[2]{ref2}
Noga Alon, Oded Goldreich, Johan Håstad, and René Peralta. Simple constructions of almost k-wise independent random variables. Random Structures \& Algorithms, 3(3):289–304, 1992.

\bibitem[3]{ref3}
Sanjeev Arora and Boaz Barak. Computational Complexity - A Modern Approach. Cambridge University Press, 2009. URL: http://www.cambridge.org/catalogue/catalogue.asp?isbn= 9780521424264.

\bibitem[4]{ref4}
Allan Borodin, Stephen Cook, and Nicholas Pippenger. Parallel computation for well-endowed rings and space-bounded probabilistic machines. Information and Control, 58(1-3):113–136, 1983.

\bibitem[5]{ref5}
Mark Braverman, Gil Cohen, and Sumegha Garg. Pseudorandom pseudo-distributions with near-optimal error for read-once branching programs. SIAM Journal on Computing, 49(5):STOC18–242–STOC18–299, 2020.

\bibitem[6]{ref6}
Eshan Chattopadhyay and Jyun-Jie Liao. Optimal error pseudodistributions for read-once branching programs. In Proceedings of the 35th Computational Complexity Conference (CCC 2020). Schloss Dagstuhl-Leibniz-Zentrum für Informatik, 2020.

\bibitem[7]{ref7}
Kuan Cheng and William M. Hoza. Hitting sets give two-sided derandomization of small space. In 35th Computational Complexity Conference (CCC 2020). Schloss Dagstuhl-Leibniz-Zentrum für Informatik, 2020.

\bibitem[8]{ref8}
Michael B. Cohen, Jonathan Kelner, John Peebles, Richard Peng, Anup B. Rao, Aaron Sidford, and Adrian Vladu. Almost linear-time algorithms for Markov chains and new spectral primitives for directed graphs. In Proceedings of the 49th Annual ACM SIGACT Symposium on Theory of Computing (STOC 2017). ACM, 2017.

\bibitem[9]{ref9}
Michael B. Cohen, Jonathan Kelner, John Peebles, Richard Peng, Aaron Sidford, and Adrian Vladu. Faster algorithms for computing the stationary distribution, simulating random walks, and more. In Proceedings of the 57th Annual IEEE Symposium on Foundations of Computer Science (FOCS 2016). IEEE, 2016.

\bibitem[10]{ref10}
Oded Goldreich. Computational complexity: a conceptual perspective. Cambridge University Press, Cambridge, 2008.

\bibitem[11]{ref11}
Oded Goldreich and Avi Wigderson. Tiny families of functions with random properties: A quality-size trade-off for hashing. Random Structures \& Algorithms, 11(4):315–343, 1997.

\bibitem[12]{ref12}
Alexander Healy and Emanuele Viola. Constant-depth circuits for arithmetic in finite fields of characteristic two. In Annual Symposium on Theoretical Aspects of Computer Science (STACS 2006). Springer, 2006.

\bibitem[13]{ref13}
William M. Hoza and David Zuckerman. Simple optimal hitting sets for small-success RL. SIAM Journal on Computing, 49(4):811–820, 2020.

\bibitem[14]{ref14}
Russel Impagliazzo, Valentine Kabanets, and Avi Wigderson. In search of an easy witness: Exponential time vs. probabilistic polynomial time. Journal of Computer and System Sciences, 65(4):672–694, 2002.

\bibitem[15]{ref15}
Russell Impagliazzo, Noam Nisan, and Avi Wigderson. Pseudorandomness for network algorithms. In Proceedings of the 26th Annual ACM SIGACT Symposium on Theory of Computing (STOC 1994). ACM, 1994.

\bibitem[16]{ref16}
Valentine Kabanets and Russell Impagliazzo. Derandomizing polynomial identity tests means proving circuit lower bounds. computational complexity, 13(1-2):1–46, 2004.

\bibitem[17]{ref17}
Noam Nisan. Pseudorandom generators for space-bounded computation. Combinatorica, 12(4):449–461, 1992.

\bibitem[18]{ref18}
Noam Nisan. RL \(\subseteq\) SC. computational complexity, 4(1):1–11, 1994.

\bibitem[19]{ref19}
Noam Nisan and Avi Wigderson. Hardness vs. randomness. Journal of Computer and System Sciences, 49(2):149–167, 1994.

\bibitem[20]{ref20}
Edward Pyne and Salil Vadhan. personal communication, February 2021.

\bibitem[21]{ref21}
Ran Raz and Omer Reingold. On recycling the randomness of states in space bounded computation. In Proceedings of the 31st Annual ACM SIGACT Symposium on Theory of Computing (STOC 1999). ACM, 1999.

\bibitem[22]{ref22}
Michael E. Saks and Shiyu Zhou. BPH SPACE(S) \(\subseteq\) DSPACE(S2/3 ). Journal of Computer and System Sceinces, 58(2):376–403, 1999.

\bibitem[23]{ref23}
Walter J. Savitch. Relationships between nondeterministic and deterministic tape complexities. Journal of Computer and System Sciences, 4(2):177–192, April 1970.
\end{thebibliography}
